\chapter*{怎样读这一卷}
\addcontentsline{toc}{chapter}{怎样读这一卷}

本卷从二二〇年的汉魏交接写到二八〇年孙吴降服。魏、蜀汉、孙吴不是三块静止的颜色：洛阳、成都与建业分别连着不同的仓储、山口、江面、户籍和边缘网络；辽东、南中与海上关系也有自己的时间和行动者。

阅读时先抓住三政权并行的年序、改元、继承与战事，再追问一项选择怎样受道路、水运、粮秣、文书和地方家庭的限制。前后相接不自动构成因果：把长期的资源条件、临近的触发、可动员的条件和可见的结果分开，才看得见战争或禅代怎样把压力转交给下一件事。

《三国志》及裴松之注、《后汉书》《晋书》和《资治通鉴》维系基本年序；《华阳国志》、走马楼吴简、石经残石、城址与区域材料补充不同尺度。魏纪年的框架、后出杂史、临湘一县的文书或使节记录，都不能单独证明全境制度、精确疆界、军队规模或行动者内心。

请把信息框当作路标而非答案：\texttt{event} 指向首次展开的事件，\texttt{turningpoint} 标示路径改变，\texttt{causalchain} 拆开条件与后果；\texttt{textwindow} 交代短引文的出处和立场；\texttt{sourcenote}、\texttt{debate} 保留材料边界与不同读法。它们应帮助读者回到连续叙事，比较魏、蜀汉、孙吴和西晋如何抵达同一场变化。
